\begin{frame}{L'annonce de la fin du livre ne date pas d'aujourd'hui...}
    \begin{block}{}
        \centering
        \small{\textit{Si par livres vous entendez parler de nos innombrables cahiers de papier imprimé, ployé, cousu, broché sous une couverture annonçant le titre de l’ouvrage, je vous avouerai franchement que je ne crois point, [...] que l’invention de Gutenberg puisse ne pas tomber plus ou moins prochainement en désuétude comme interprète de nos productions intellectuelles.
        Je crois donc au succès de tout ce qui flattera et entretiendra la paresse et l’égoïsme de l’homme ; l’ascenseur a tué les ascensions dans les maisons ; le phonographe détruira probablement l’imprimerie.}\\
        \textbf{"La fin des livres", Contes pour bibliophiles, Octave Uzanne (illustr. Albert Robida), 1894.}}
    \end{block}
    = littérature d'anticipation.
\end{frame}

\begin{frame}{L'annonce de la fin du livre ne date pas d'aujourd'hui...}
    \begin{block}{}
        \centering
        \small{\textit{Pour le livre, ou disons mieux, car alors les livres auront vécu, pour le novel ou storyographe, l’auteur deviendra son propre éditeur, afin d’éviter les imitations et contrefaçons ; il devra préalablement se rendre au Patent Office pour y déposer sa voix et en signer les notes basses et hautes, en donnant des contre-auditions nécessaires pour assurer les doubles de sa consignation. Aussitôt cette mise en règle avec la loi, l’auteur parlera son œuvre et la clichera sur des rouleaux enregistreurs et mettra en vente lui-même ses cylindres patentés, qui seront livrés sous enveloppe à la consommation des auditeurs.}\\
        \textbf{"La fin des livres", Contes pour bibliophiles, Octave Uzanne (illustr. Albert Robida), 1894.}}
    \end{block}

En résumé:
\begin{itemize}
    \item L'auteur s'auto-éditera.
    \item On laissera tomber l'écriture pour la voix (podcast)
\end{itemize}
\end{frame}